% Generated by tools/edn_latex.py from reports/edn.md.
% Do not edit: change reports/edn.md and run `make edn`.
\ednentry{
  id = {EDN-23},
  title = {Read the condition flags as one-sided assertions instead of encoding them three-valued},
  date = {2026-09-29},
  milestone = {M3: Quality Assurance},
  topic = {Feature Encoding and Data Validation},
  participants = {@lukas2510},
  decision = {\texttt{has\_\allowbreak{}full\_\allowbreak{}service\_\allowbreak{}history}, \texttt{non\_\allowbreak{}smoking} and \texttt{is\_\allowbreak{}rental} stay plain booleans and are read as ``the seller asserted this'' versus ``the seller did not assert this'', never as ``yes'' versus ``no''. No three-valued (true / false / unknown) encoding is built and no \texttt{False} is mapped to missing. \texttt{is\_\allowbreak{}used} is excluded entirely because \texttt{offer\_\allowbreak{}type} already carries the same information reliably. The Great Expectations suite only checks that the columns are boolean and non-null.},
  rationale = {\emph{Alternatives considered:}
\begin{itemize}[leftmargin=*, topsep=2pt]
\item \textbf{Option A (chosen): keep them binary, fix the interpretation.}
\newline Pros: no pipeline, API or serialisation complexity; loses nothing, because the ambiguity is in the source and no encoding can resolve it; the risk it addresses (somebody reading a \texttt{False} as a denial) is a documentation and naming problem, so it is solved where it actually lives, in the problem spec, the dataset card and the model card.
\newline Cons: the column name still reads like a two-sided statement, so the documentation has to carry the caveat; a reader who skips it can still misread the feature.
\item \textbf{Option B: encode three-valued, or map \texttt{False} to missing.}
\newline Pros: makes the ambiguity visible in the data itself rather than only in prose; the shape a careful reviewer expects.
\newline Cons: it is not actually three-valued, because we cannot tell a real \texttt{False} from an unknown, so the third state would be empty and the encoding would misrepresent the problem as solved; for a gradient-boosted tree a column holding only \texttt{\{True, NaN\}} splits exactly like \texttt{\{True, False\}}, so the model gains no information at all; adds an extra state to the API schema, the training pipeline and the expectations for that zero gain.
\item \textbf{Option C: drop all affected flags.}
\newline Pros: the easiest position to defend; no ambiguous column ever reaches the model.
\newline Cons: discards the half of the signal that is reliable, since a \texttt{True} is a real assertion; full service history and non-smoking are exactly the kind of attributes that move a used-car price.
\item \textbf{Option D: recover the true value from \texttt{description} or the equipment lists.}
\newline Pros: would turn the flags into genuine two-sided features.
\newline Cons: needs multilingual free-text parsing, which the problem spec puts out of scope; produces labels with no ground truth to validate against.
\end{itemize}
\par
\emph{Why:} The flags are one-sided by construction: a listing form offers a checkbox, so a tick is evidence and an empty box is the absence of evidence. 14,744 rows flagged as neither used, new nor pre-registered, and 18,108 scoped rows with \texttt{is\_\allowbreak{}used =\allowbreak{} False} while \texttt{offer\_\allowbreak{}type =\allowbreak{} U}, are the direct measurement of that. Since no information exists to separate ``no'' from ``not stated'', every encoding choice carries identical information, and the one that costs nothing wins. This deliberately overrides the original proposal in issue \#25 to encode the flags three-valued.},
  ai = {Information seeking, Alternative generation, Alternative assessment, Recommendation},
  response = {Accepted},
  assessment = {The three-valued encoding was AI's own earlier proposal, written into issue \#25. On re-reading the issue, AI argued against it: it showed that \texttt{\{True, NaN\}} and \texttt{\{True, False\}} are informationally identical for the tree models we use, so the encoding would add complexity without any effect on the model, and recommended fixing the interpretation instead. Lukas accepted the reversal because the argument was concrete rather than stylistic.},
  aievidence = {Claude Code session on 2026-09-29 analysing \href{https://github.com/mlops-2627q1-mds-upc/MLOPS_Recommenditos/issues/25}{issue \#25}: AI presented the options as a decision question, Lukas chose ``keep them binary''.},
  evidence = {\href{https://github.com/mlops-2627q1-mds-upc/MLOPS_Recommenditos/issues/25}{Issue \#25}, \href{https://github.com/mlops-2627q1-mds-upc/MLOPS_Recommenditos/pull/24}{PR \#24}, \href{https://github.com/mlops-2627q1-mds-upc/MLOPS_Recommenditos/blob/main/docs/docs/problem-spec.md}{problem specification} section 4.},
}
